\subsection{The common quotient with the degree-eight action}

\begin{statement}{D.38}{Proved}{A simultaneous mixed Fermat and degree-eight shadow}
Retain $L,L_4,J_{12}$ from D.37 and the marked free
rank-four group $K$ from D.34. Put
\[
 E_2=B/[L_4,L_4]L_4^4,\qquad T_3=B/D_9(K).
                                                               \tag{D.38.1}
\]
The maximal common marked quotient of $J_{12}$ and $T_3$,
with their displayed maps from $B$, is
\[
 C=B/\bigl([L_4,L_4]L_4^4D_9(K)\bigr).
                                                               \tag{D.38.2}
\]
The image of $D_9(K)$ in $E_2$ lies in $L_4^{\rm ab}/4$
and has exact group type $C_4^{51}\oplus C_2^{265}$.
Consequently
\[
 \ker(C\to H_4)\simeq C_4^{69}\oplus C_2^{265},
 \qquad |C|=2304\,2^{403}=18\,2^{410}.                    \tag{D.38.3}
\]

There is an actual generating finite GT0 shadow on
\[
 \mathcal J
 =B/\bigl([L,L]L^{12}\cap D_9(K)\bigr)
 \simeq J_{12}\mathbin{\times_C}T_3.                      \tag{D.38.4}
\]
Its multiplier is one. Before the common Ihara gauge,
its $T_3$ component is exactly the recorded D.33 action,
and its two conjugators have exactly the prescribed
common images in $C$. Its new $J_{12}$ component retains
D.37's normalized $G_6$ and $J_3$ actions. It replaces
the recorded D.37 $J_{12}$ action: both normalized
generator images differ already modulo two.

The common gauge exponent is $k=-774$. Since
$x^k=c_0^{258}$, the normalized $T_3$ component is
$\operatorname{Int}(c_0^{258})\alpha_{33}$. It therefore
has exactly the same image in D.35's quotient by
$\langle\operatorname{Int}(c_0)\rangle$. The compact
evaluations, in D.36's bases, are
\[
 \rho_{C_3,16}=
 \begin{pmatrix}3&8\\8&11\end{pmatrix},
 \qquad \rho_{C_4,3}=I.                                  \tag{D.38.5}
\]
The cubic reduction modulo four is $-I$, so the same
arithmetic-graph violation and the same extra fixed-fibre
pair $T_4=I,\ W_3=[i]$ persist.

The orders and projection kernels are
\[
 \begin{gathered}
 |\mathcal J|=2^{13472}3^{1156},\\
 |\ker(\mathcal J\to J_{12})|=2^{11158},\qquad
 |\ker(\mathcal J\to T_3)|=2^{1903}3^{1154}.
 \end{gathered}                                           \tag{D.38.6}
\]
Its cusp still has order 288 and its finite-shadow
$N_{\rm ord}$ is 144. For the two exact source maps
\[
 \widehat B\twoheadrightarrow\mathcal J,\qquad
 \widehat F\twoheadrightarrow\ker(\mathcal J\to S_3),
\]
the kernel is the closure of the intersection in
(D.38.4), taken in the respective source. It contains the explicit
nonzero word
\[
 [a_1^{12},a_2^{12}]^2,                                  \tag{D.38.7}
\]
where $a_1,a_2$ are the first two marked generators of $K$.
This word has a nonzero degree-sixteen mod-two Magnus
term. Thus the simultaneous finite action leaves an
explicit nontrivial deck kernel for the full lifting
test. D.45 and its source-transport supplement exclude a full
profinite lift of this particular action, whose degree-eight
component differs from D.33 by the stated source cusp-inner map.

This addresses Part D(e), section 51.16, folio 734 verso,
scan 174, and D-after-LM2(1), section 51.15, folios
725--727, scans 156--160. The comparison retains D.33's
degree-eight component and the actual two cover maps of
D.37; it does not assert compatibility with D.33's entire
infinite class-two system.
\end{statement}
\begin{proof}
We compute the common quotient before solving any
conjugator equation. Let
$V=L_4^{\rm ab}/4\simeq(\mathbb Z/4)^{385}$, the normal
coefficient subgroup of $E_2$, and let $G$ be the image
of $K$ there. The quotient $A_0=K/L_4$ has order 128.
It embeds in
\[
 O_2(Q)\times(\mathbb Z/4)^2,
                                                               \tag{D.38.8}
\]
because the affine linear part of every $K$ element is
the identity. The two factors have class at most two
and exponent four. Thus $A_0$ has those properties,
and its commutator group has exponent two. The program
also verifies these statements on every pair of its
128 elements.

It follows that $\gamma_3G\subset V$, every element of
$G$ has order dividing 16, and every element of
$\gamma_2G$ has order dividing 8. The latter assertion
uses $(\gamma_2G)^2\subset V$. Write $U=\gamma_3G$,
and let $I$ be the augmentation ideal for the action
of the four actual marked $K$ generators on $V$.
Since $V$ is abelian, for every $j\ge3$,
\[
 \gamma_jG=I^{j-3}U.
                                                               \tag{D.38.9}
\]
The Zassenhaus product
$D_n(G)=\prod_{i2^r\ge n}(\gamma_iG)^{2^r}$ therefore
gives the exact equality
\[
 W:=\operatorname{im}\bigl(D_9K\to E_2\bigr)
   =\gamma_9G+2\gamma_5G=I^6U+2I^2U.                     \tag{D.38.10}
\]
Indeed $G^{16}$, $(\gamma_2G)^8$ and
$(\gamma_3G)^4$ vanish, and the remaining terms are
contained in the two displayed summands. The equality
is functorial under the surjection $K\twoheadrightarrow G$,
since every lower-central commutator and every power
in the Zassenhaus product has a preimage of the same form.

All groups in (D.38.10) are computed in one integral
graph module. The $H_4$ graph has 2304 edges and 1920
vertices; its 385 non-tree edges give the identity
minor of its cycle matrix. The four marked $K$ words
are evaluated in the faithful model
$(\mathbb Z/4)[H_4]\rtimes H_4$. Their twenty-four
basic triple commutators lie in $V$. Taking their
module span and repeatedly adjoining their four
translates gives their full normal closure $U$.
This is $\gamma_3G$: making the basic generator
commutators central is precisely the quotient of
class at most two.

The successive exact group types are
\[
 \begin{array}{c|rr}
 j&\text{number of }C_4\text{ factors in }\gamma_jG&
       \text{number of }C_2\text{ factors}\\\hline
 3&375&1\\
 4&355&11\\
 5&316&38\\
 6&252&85\\
 7&178&134\\
 8&106&146\\
 9&51&127
 \end{array}                                               \tag{D.38.11}
\]
Adjoining $2\gamma_5G$ to the last row gives
$W\simeq C_4^{51}\oplus C_2^{265}$.
The calculation performs elimination over $\mathbb Z/4$:
odd pivots give independent order-four generators,
and the remaining even rows reduce to exact binary
elimination. It also checks membership of every
generator translate, including the full $s,r$ action.
This verifies the normal closure and every subsequent
module product, rather than only their rational ranks.

Quotienting the free coefficient module $V$ by $W$
therefore leaves
$C_4^{385-51-265}\oplus C_2^{265}$, proving (D.38.3).
The universal relative pro-two quotient of $J_{12}$
over $Q$ is $E_2$, by D.37. On the other hand,
$D_9K\subset N$, and $T_3\to Q$ has a two-group kernel.
Every common quotient of $J_{12}$ and $T_3$ over $Q$
must consequently factor through $E_2$. The largest
common quotient of $E_2$ and $B/D_9K$ is exactly
$E_2/W$, by the definition of the product of the two
normal kernels. This proves (D.38.2), in both
directions, and the normal-subgroup Chinese remainder
isomorphism proves (D.38.4).

We now impose the exact common images of the recorded
D.33 conjugators. Both long words are reconstructed
from its retained textual coefficients, without
rerunning or changing any earlier solver. Their
SHA256 hashes agree with the original reconstruction
certificate, and their full cusp word has the
recorded length 12921196. Evaluation in $E_2$ gives
paths $P_A^*$ and $P_\beta^*$. Their cusp defect lies
in $W$, providing a direct check that the recorded
$T_3$ action gives its stated common-quotient action.

The projection $H\to H_4$ sends each graph edge to
its projected edge. On the cycle bases it gives an
exact matrix
\[
 S_4:\mathbb Z^{1153}\longrightarrow\mathbb Z^{385}.
                                                               \tag{D.38.12}
\]
Let $a,b$ be the two unknown $L$ cycle corrections,
and let $Z$ be the $385\times316$ matrix of the
recorded generators of $W$. In addition to D.37's
complete cusp equation modulo four, impose
\[
 \begin{split}
 S_4a-Zu_A&=(P_A^*-\pi P_A)_{\rm cycles},\\
 S_4b-Zu_\beta&=(P_\beta^*-\pi P_\beta)_{\rm cycles}.
 \end{split}                                             \tag{D.38.13}
\]
These equations state equality of the actual
conjugators in $C$, including all central coordinates.
The prior exact $G_2$ conjugator constraints are included
as twice their binary equations. There are 2117
equations and 2938 unknowns over $\mathbb Z/4$.

Unit-pivot elimination has rank 1588. The residual even
system, divided by two, has binary rank 263. The full
solution space is an explicitly parametrized torsor
under
\[
 C_4^{1087}\oplus C_2^{263}.                              \tag{D.38.14}
\]
The first recorded solution has a full rank 1153 map
on $L^{\rm ab}/2$. It is combined with exactly the
same mod-three correction as D.37. Thus it gives an
automorphism of $J_{12}$: its mod-four kernel matrix
is invertible because its mod-two reduction is,
and its unchanged mod-three matrix was already
verified to be invertible.

The resulting 669040-letter cusp word is evaluated
in all 6912 coordinates modulo12. Every one of the
1153 free-kernel basis images is independently
computed by generator substitution and agrees with
the graph map. The program also checks (D.38.13)
against the stored $W$ coefficients and against
the separate exact membership reduction.

The two component automorphisms agree on $C$ because
their conjugators do. More explicitly, their
compositions from $B$ to $C$ agree on $s,r$; one
kills the $J_{12}$ kernel and the other kills
$D_9K$, so this common homomorphism factors through
their product kernel. It is onto because either
component is an automorphism. The product of the
two component automorphisms therefore restricts
to an automorphism of the fibre product in
(D.38.4). This is generation for the actual joint
finite group, not a rank assertion for a linear
approximation.

Both component conjugators lie over the identity
permutation, so their compatible pairs lift to
actual words of $F$ under the surjection
$B\to\mathcal J$. For the recorded $J_{12}$
representative,
$\beta^{\rm ab}=666x-108y$. The gauge exponent is
$-108-666=-774$. Every other $F$ lift with the same
$J_{12}$ image changes this abelianization by an
element of $144\mathbb Z^2$: the kernel is
$[L,L]L^{12}$ and $L\subset F_{12}$. Write the
abelianization of a compatible lift as
$(666+144u,-108+144v)$. With the displayed gauge
$k=-774$, the word
$f=s(x^k\beta)s(x^k\beta)^{-1}$ has abelianization
\[
 f^{\rm ab}=(144(v-u),144(u-v))\equiv(0,0)\pmod{144}.
\]
Both $x$ and $y=sxs$ have order 144 in $\mathcal J$:
$x$ has orders 144 and 48 in the two components.
Putting $\epsilon=v-u$, the word
\[
 f_0=x^{-144\epsilon}y^{144\epsilon}f\in F'
\]
has the same joint finite image as $f$. Thus the
joint $f$ coset contains an actual commutator word,
including the effect of every compatible $T_3$ lift.

The two $m=0$ hexagons follow as in D.25 from the
torsion and cusp equations, and are also checked in
the new faithful $J_{12}$ component. They hold in
$T_3$ by the same identities for its recorded action
and the common inner gauge. Together with the
joint automorphism this verifies the finite GT0-shadow
conditions of \cite[Definitions 2.1 and 2.3]{dBDHP2026}.
The preimage of the joint kernel under $B_3\to B$
is normal, has finite index, and lies in $PB_3$, since
the joint quotient retains $B\to S_3$. The braid
centre maps to one. The cusp has order
$\operatorname{lcm}(288,96)=288$, and both pure loops
have order 144; hence the pulled-back kernel has
$N_{\rm ord}=144$ and the multiplier is one.

The gauge is particularly explicit on the degree-eight
component:
$x^{-774}=t^{1548}=(t^6)^{258}=c_0^{258}$.
It therefore changes the recorded $T_3$ automorphism
by exactly the inner automorphism stated in the
theorem. It fixes its class under D.35's specified
cusp-inner quotient. On the other component,
the normalized new and old images of each of $s,r$
have nonzero differences on 5188 edge coordinates
modulo12, and already differ modulo two. Their
projections to $G_6$ agree, and their reductions
modulo three on the full $H$ module agree. These
checks prove the asserted replacement and both
retained normalized actions.

The compact-cover evaluation is repeated using the
new normalized $97388$-letter $F$ word. All
Schreier-generator substitutions for $C_3,C_4,C_{12}$
and both compact projection identities pass. They
again give $-I$ on cubic four-torsion and $I$ on
quartic three-torsion.

For the finer cubic component, $D_9K$ maps to
$16K^{\rm ab}$, so $T_3$ retains the full
$K^{\rm ab}/16$ action. The exact compact change of
basis from D.34 to D.36 is
\[
 B_{\rm c}=\begin{pmatrix}-1&1\\-1&0\end{pmatrix}.
                                                               \tag{D.38.15}
\]
It is obtained by rewriting the four marked $K$
words in the cyclic-three Schreier basis and
applying the two integral filling maps. Its
determinant is one, it intertwines the marked
deck matrices, and
$B_{\rm c}M_{16}B_{\rm c}^{-1}=M_{16}$ for
$M_{16}=\left(\begin{smallmatrix}3&8\\8&11\end{smallmatrix}\right)$.
The gauge is a $K$-inner cusp conjugation, so it acts
trivially on this homology. This proves (D.38.5)
through the exact comparison maps. D.36's arithmetic
graph therefore still excludes a Galois realization;
the unchanged $G_6$ CM pair gives the independent
exclusion of D.25.

For the orders, the restricted free-Lie dimensions
of rank four in degrees one through eight are
\[
 4,\ 10,\ 20,\ 70,\ 204,\ 690,\ 2340,\ 8230.
\]
Their sum is 11568, giving $|T_3|=18\,2^{11568}$.
Equations (D.37.2), (D.38.3) and the fibre-product
formula give every order in (D.38.6).

Both kernels in (D.38.4) are normal and have finite
index in $B$, and their intersection is contained in $F$.
Consequently their completions give exactly the two
source maps in the statement; the open embedding
$\widehat F\hookrightarrow\widehat B$ identifies their
kernels. Finally, each marked $a_i^{12}$ lies in $L$.
Indeed the $K$ image in $H$ embeds in
$O_2(Q)\times(\mathbb Z/12)^2$ and has exponent
dividing 12; the two specific words are also
checked directly. Thus (D.38.7) lies in $[L,L]$.
In the free mod-two Magnus algebra, the first term
of $a_i^{12}-1$ is $X_i^4$. The first term of the
commutator is consequently
$X_1^4X_2^4+X_2^4X_1^4$, and squaring gives four
distinct degree-sixteen monomials. It is nonzero
and belongs to $D_{16}K\subset D_9K$. This proves
both its inclusion in the stated remaining kernel
and its nontriviality in the full free pro-two,
hence full profinite, source. The D.45 obstruction, transported
through the exact cusp-inner gauge above, excludes a full profinite
lift of this recorded finite shadow. The inverse-system test of
D.22 remains required for a replacement realizing its CM pair.
\end{proof}
